\section{Introducción}

El presente capítulo documenta la materialización del diseño especificado en el
capítulo anterior, así como la verificación experimental de su desempeño. Su
propósito es doble: por un lado, describir la implementación efectiva de cada
subsistema ---el motor de percepción, la plataforma de gestión y la interfaz de
usuario--- junto con las decisiones técnicas adoptadas durante la construcción;
por otro, presentar los resultados obtenidos en las pruebas realizadas bajo
condiciones representativas del entorno de estudio.

A diferencia del capítulo de diseño, que establece la arquitectura conceptual y
los criterios de aceptación, este capítulo se apoya en artefactos verificables:
fragmentos de código, esquemas de base de datos, archivos de configuración,
métricas medidas y capturas de la aplicación en funcionamiento. La trazabilidad
entre ambas fases se mantiene mediante los identificadores de requerimiento
definidos en las Tablas \ref{tab:requerimientos_funcionales} y
\ref{tab:requerimientos_no_funcionales}, de modo que cada requisito pueda
vincularse con el componente que lo satisface y con la prueba que lo verifica.

El capítulo se organiza siguiendo la secuencia de construcción del sistema. En
primer lugar, se describe el modelo de implementación: el entorno tecnológico,
la arquitectura de despliegue y la organización del repositorio. A continuación,
se documenta la construcción del conjunto de datos, etapa habilitante para el
entrenamiento de los modelos. Posteriormente, se detalla la implementación del
subsistema de visión artificial y del subsistema de gestión. Luego se presenta
la estrategia de optimización y despliegue en hardware de borde. Finalmente, se
exponen el protocolo experimental y los resultados obtenidos, y se muestran las
interfaces del sistema en operación.
